\chapter{Manual de instalación y despliegue}
\label{app:manual_instalacion}

\ph{Este apéndice describe los pasos para instalar y desplegar el sistema en
un entorno limpio. Debe ser lo bastante detallado para que un tercero
(p.\,ej.\ otro estudiante o un evaluador) pueda reproducir el despliegue
sin contacto con el autor/a.}

\section{Requisitos previos}
\label{sec:requisitos_previos}

\ph{Lista del software, hardware y credenciales necesarios antes de empezar.
Indica versiones concretas siempre que sea posible.}

\begin{itemize}
  \item \textbf{Sistema operativo:} \ph{p.\,ej.\ Linux Ubuntu 24.04 LTS / macOS / Windows 11.}
  \item \textbf{Lenguaje y runtime:} \ph{p.\,ej.\ Python 3.x / Node 20 / Java 17.}
  \item \textbf{Gestor de paquetes:} \ph{pip / npm / maven\dots}
  \item \textbf{Servicios externos:} \ph{base de datos, broker de mensajería,
        almacenamiento\dots}
  \item \textbf{Otras herramientas:} \ph{Docker, Git, etc.}
\end{itemize}

\section{Obtención del código fuente}
\label{sec:obtencion_codigo}

\ph{Indica el repositorio donde está alojado el código y cómo clonarlo:}

\begin{verbatim}
git clone <URL del repositorio>
cd <nombre del proyecto>
\end{verbatim}

\section{Instalación de dependencias}
\label{sec:instalacion_dependencias}

\ph{Comandos para instalar las dependencias del proyecto en un entorno
aislado (virtualenv, contenedor, etc.). Si hay variables de entorno o
ficheros de configuración que el usuario debe rellenar, lístalos aquí.}

\begin{verbatim}
[Comandos de instalación de dependencias]

[Variables de entorno requeridas:
 - VAR_1 = ...
 - VAR_2 = ...]
\end{verbatim}

\section{Configuración del entorno}
\label{sec:configuracion_entorno}

\ph{Pasos adicionales antes de poder ejecutar el sistema: creación de la base
de datos, ejecución de migraciones, carga de datos iniciales, generación de
claves, configuración del proxy inverso, etc.}

\begin{verbatim}
[Pasos de configuración inicial]
\end{verbatim}

\section{Ejecución}
\label{sec:ejecucion}

\ph{Comando o comandos para arrancar el sistema en local. Indica también
en qué URL/puerto queda disponible.}

\begin{verbatim}
[Comando de arranque]
\end{verbatim}

\section{Despliegue en producción}
\label{sec:despliegue_produccion}

\ph{Si el sistema se despliega en un entorno productivo (servidor propio,
nube, PaaS\dots), describe el proceso paso a paso: construcción de la
imagen, publicación en el registro, despliegue en el host, configuración
del proxy/HTTPS, etc.}

\begin{verbatim}
[Pasos de despliegue en producción]
\end{verbatim}

\section{Verificación de la instalación}
\label{sec:verificacion}

\ph{Cómo comprobar que la instalación ha sido correcta: ejecución de la
suite de pruebas, acceso a una pantalla concreta, llamada a un endpoint
de \textit{health check}, etc.}
